\documentclass[11pt,letterpaper]{article}
\usepackage[T1]{fontenc}
\usepackage{lmodern}
\usepackage{microtype}
\usepackage[margin=1in]{geometry}
\linespread{1.0}
\setlength{\parskip}{0pt}
\setlength{\parindent}{1.5em}
\usepackage{amsmath,amssymb,amsthm}
\usepackage{graphicx}
\usepackage{xcolor}
\usepackage{tikz}
\usepackage{booktabs,tabularx,array,longtable}
\newcolumntype{L}[1]{>{\raggedright\arraybackslash}p{#1}}
\usepackage{listings}
\lstset{basicstyle=\ttfamily\footnotesize,breaklines=true,columns=fullflexible,frame=single,framesep=4pt,showstringspaces=false,captionpos=b}
\usepackage[hidelinks]{hyperref}
\hypersetup{pdftitle={Obstructions to Adelic Polarization},
  pdfauthor={Matthew Long; The YonedaAI Collaboration}}
\usepackage{cleveref}
\usepackage{everypage}
\newtheorem{theorem}{Theorem}[section]
\newtheorem{proposition}[theorem]{Proposition}
\newtheorem{lemma}[theorem]{Lemma}
\newtheorem{corollary}[theorem]{Corollary}
\theoremstyle{definition}
\newtheorem{definition}[theorem]{Definition}
\theoremstyle{remark}
\newtheorem{remark}[theorem]{Remark}
\newcommand{\R}{\mathbb R}
\newcommand{\C}{\mathbb C}
\newcommand{\Z}{\mathbb Z}
\newcommand{\T}{\mathcal T}
\newcommand{\B}{\mathcal B}
\newcommand{\supp}{\operatorname{supp}}
\newcommand{\ord}{\operatorname{ord}}
\newcommand{\Id}{\operatorname{Id}}
\newcommand{\Span}{\operatorname{span}}
\newcommand{\Dom}{\operatorname{Dom}}
\definecolor{grokgray}{RGB}{110,110,110}
\AddEverypageHook{%
  \ifnum\value{page}=1
    \begin{tikzpicture}[remember picture,overlay]
      \node[rotate=90,anchor=south,font=\footnotesize\sffamily\color{grokgray},inner sep=0pt]
      at ([xshift=30pt,yshift=0.5\paperheight]current page.south west)
      {GrokRxiv:2026.09.polarization-obstructions\quad [\,math.NT\,]\quad 08 Sep 2026};
    \end{tikzpicture}
  \fi
}
\title{Obstructions to Adelic Polarization\\[4pt]
{\normalsize Part II of \emph{Adelic Polarization Conjecture}}}
\author{Matthew Long\\
\small The YonedaAI Collaboration, YonedaAI Research Collective\\
\small Chicago, IL\\
\small\texttt{matthew@yonedaai.com}}
\date{September 2026}

\begin{document}
\maketitle
\begin{abstract}
Arithmetic polarization through semilocal adelic spaces requires prime transitions that preserve arithmetic comparison as well as a geometric sign. We examine two particular ways of imposing these requirements. The ambient polar normalization of a known weighted entire-function carrier divides boundary values by the modulus of a local Euler factor. We prove that this operation cannot send any nonzero entire function to an entire function. The proof uses a simple zero of the squared Euler modulus and the parity of vanishing orders. The intrinsic polar decomposition on a closed Sonin subspace has a compressed metric operator and is not excluded by this argument. We derive its adjoint and explain the additional test-map and symmetry identities it must satisfy. A separate obstruction comes from exact compact smooth tests whose two pole moments vanish. At the arithmetic prime 11, three and four translated bumps give local quadratic values proportional to $-46/11$ and $(28\sqrt{11}-92)/11$. These opposite signs rule out an orthogonal nonpositive enlargement with the prescribed positive local Weil increment on all intermediate tests. A finite packet calculation extends the sign result to every prime. The obstruction does not apply to sign imposed only at support-complete stages, or to correspondences with mixed pairings and explicit correction terms. The resulting formulation retains those terms and leaves the geometric sign and full arithmetic comparison as construction problems.
\end{abstract}

\section{Introduction}
\label{sec:introduction}

A topological identification of semilocal Sonin spaces does not identify
their inner products. This distinction matters when those spaces are used
as a possible analytic realization of arithmetic cohomology. The
construction of Connes, Consani and Moscovici provides a common space of
entire functions with norms that change when a prime is adjoined
\cite[Sections 4.7 and 4.8]{CCM24}. An arithmetic polarization would need
more than equivalent norms: its pairing must arise from arithmetic
duality and reproduce the signed local terms of the explicit formula.

For a prime $p$, put $L=\log p$ and $r=p^{-1/2}$. In the common carrier,
the finite change of weight is multiplication by
\begin{equation}
 \omega_p(t)=|1-re^{-itL}|^2=1+r^2-2r\cos(tL).
 \label{eq:weight}
\end{equation}
This positive function lies between $(1-r)^2$ and $(1+r)^2$.
The multiplier $\omega_p-1$ is denoted $d_p$ in Part I.
The positivity of $\omega_p$ gives bounded comparison of the two norms. It does not
give a sign to their difference, and it does not give a sign to the
logarithmic derivative that produces the local Weil distribution.
Part I establishes these local analytic identities with a fixed Fourier
convention \cite{PartI}.

One possible repair is to replace the carrier identity by its polar
isometry. In the ambient weighted $L^2$ spaces this means division by
$\sqrt{\omega_p}$. The operation is well-defined, bounded and invertible
on boundary functions. Theorem~\ref{thm:polarization-obstructions:entire}
shows that it destroys the entire-function condition for every nonzero
input. Squaring a putative boundary identity produces an identity of
entire functions. Its two sides have incompatible orders at zero.
This obstruction does not depend on the growth or reproducing kernel
of the Sonin space.

Polar decomposition within a closed Sonin space is a different
operation. The adjoint of the restricted carrier identity includes the
orthogonal projection onto that space. The associated positive operator
is a compression of multiplication by $\omega_p$, and its inverse square
root need not be pointwise multiplication. We calculate this operator
and its transformed projection in
Proposition~\ref{prop:polarization-obstructions:compression}.
The calculation preserves a possible intrinsic normalization while
making its required arithmetic compatibility precise.

A second proposed repair is geometric: keep the old primitive sector
isometrically and add an orthogonal nonpositive sector at each prime.
The new sector then contributes a nonpositive quadratic increment.
The required intersection increment, with the convention
$B_W=-I$, is $+W_p$. That increment has both signs, even after the
two pole moments vanish. We give exact smooth witnesses at $p=11$,
then a finite packet proof for every prime. The proof evaluates the
actual distribution on autocorrelations; a finite matrix computation
alone would not establish this statement about smooth tests.

The resulting exclusion has a specific support hypothesis. A new prime
contributes zero to tests whose support was already complete for the
old finite place set. The positive witness therefore obstructs the
orthogonal sign ansatz on a nontrivial intermediate transition. It
does not contradict a construction requiring its final sign only
after all primes relevant to the support have been included.

Neither obstruction establishes a sign for the full Weil form. They
reject two specified transitions and identify conditions a replacement
must meet. The entire carrier, the local explicit formula and the
cohomological motivation all precede this work. We make no claim that
these elementary obstructions have priority over every treatment of
Euler factors or compressed operators.

\section{Mathematical framework}
\label{sec:framework}

The conventions of Part I, Definition 2.1, are used throughout.
Thus $\T=C_c^\infty(\R,\C)$, convolution uses Lebesgue measure,
$f^*(x)=\overline{f(-x)}$, and the centered Laplace transform is
$F_f(z)=\int_\R f(x)e^{zx}\,dx$. The Fourier transform is
$\widehat f(t)=F_f(-it)$. These are imported conventions; the
geometric test correspondence considered later is not part of their
definition.

Part I, Definition 3.1, fixes
$W=P-A-\sum_pW_p$. In particular the local distribution is
\begin{equation}
 W_p(h)=L\sum_{k\geq1}r^k\{h(kL)+h(-kL)\},
 \qquad h\in\T.
 \label{eq:local-Weil}
\end{equation}
The sum is finite on each compactly supported test. The corresponding
Hermitian forms are $B_p(f,g)=W_p(f*g^*)$ and $B_W(f,g)=W(f*g^*)$.
The minus sign in the full
Weil distribution follows from the classical explicit formula with
its archimedean subtraction \cite[Appendix B]{CC20}. Consequently an
intersection form intended to satisfy $B_W=-I$ has local increment
$+B_p$, provided no further term changes at the same transition.

The subspace of tests with the two pole moments zero is
\begin{equation}
 \T_0=\{f\in\T:F_f(1/2)=F_f(-1/2)=0\}.
 \label{eq:pole-tests}
\end{equation}
No reality or evenness condition is imposed on $\T_0$. Our explicit
witnesses will be real, which suffices to prove both signs on this
complex space. The pole term vanishes on $f*g^*$ when $f,g\in\T_0$;
its vanishing does not remove the archimedean term or any finite prime.

For the analytic carrier, write $F_t=F(1/2+it)$ and
\begin{equation}
 q_S(F,G)=\int_\R F_t\overline{G_t}w_S(t)\,dt,
 \qquad w_{S\cup\{p\}}=\omega_p w_S.
 \label{eq:carrier-norms}
\end{equation}
These are the common-carrier norms imported from Part I,
Section 8, equations (46) and (47), and from equations (59)
and (62) of \cite{CCM24}.
The positive analytic form $q_S$ and the hypothetical geometric
form $I_S$ are distinct objects. Equating either one with a desired
arithmetic form is a comparison theorem to prove, not a convention.

All Hermitian forms in this paper are linear in the first argument.
If $A:H_0\to H_1$ is bounded, its adjoint satisfies
$\langle Ax,y\rangle_1=\langle x,A^*y\rangle_0$.
Whenever an orthogonal projection onto an entire-function space is
used below, that space is assumed to be a closed subspace of the
ambient Hilbert realization. This hypothesis is essential for the
projection calculation. The entire obstruction itself needs no
Hilbert-space hypothesis.

\section{Entire normalization}
\label{sec:entire}

\begin{definition}
\label{def:polarization-obstructions:discriminant}
For a real number $p>1$, the entire Euler denominator is
\begin{equation}
 \Delta_p(z)=(1-e^{-Lz})(1-e^{L(z-1)}),\qquad L=\log p.
 \label{eq:discriminant}
\end{equation}
The powers $p^z$ always mean $e^{z\log p}$, where $\log p$
is the unique real number $L$ satisfying $e^L=p$.
\end{definition}

For $z=1/2+it$, the two factors in \eqref{eq:discriminant} are
complex conjugates. Hence $\Delta_p(1/2+it)=\omega_p(t)>0$. Away from the
critical line the same function has zeros, and their orders constrain
any holomorphic square root.

\begin{lemma}
\label{lem:simple-zero}
The zeros of $\Delta_p$ are simple and occur at
\begin{equation}
 z=\frac{2\pi i k}{L}\quad\hbox{or}\quad
 z=1+\frac{2\pi i k}{L},\qquad k\in\Z.
 \label{eq:zero-lattices}
\end{equation}
In particular $\Delta_p(0)=0$ and $\Delta_p'(0)=L(1-p^{-1})\ne0$.
\end{lemma}
\begin{proof}
The equation $e^{-Lz}=1$ is equivalent to the first lattice in
\eqref{eq:zero-lattices}; changing the sign of $k$ does not change
the set. Its factor has derivative $Le^{-Lz}=L$ at every zero.
At such a point the other factor is $1-p^{-1}$.
The second factor vanishes at the second lattice and has derivative
$-Le^{L(z-1)}=-L$ there. The first factor is then $1-p^{-1}$.
The two lattices have different real parts and cannot intersect.
Thus exactly one factor vanishes at each zero, with nonzero
derivative, which proves the claim.
\end{proof}

\begin{lemma}
\label{lem:boundary-identity}
Let $F,G$ be entire functions. If
\begin{equation}
 G(1/2+it)=\omega_p(t)^{-1/2}F(1/2+it)
 \label{eq:boundary-normalization}
\end{equation}
for Lebesgue-almost every $t\in\R$, then
\begin{equation}
 \Delta_p(z)G(z)^2=F(z)^2\qquad(z\in\C).
 \label{eq:entire-identity}
\end{equation}
The same conclusion holds if the exceptional set is null for a
measure $w(t)\,dt$ with $w>0$ almost everywhere.
\end{lemma}
\begin{proof}
The difference of the two sides of
\eqref{eq:boundary-normalization} is continuous in $t$, because
$\omega_p(t)$ is continuous and bounded below by $(1-r)^2>0$.
A continuous function that is zero almost everywhere on $\R$ is
zero everywhere: if its value were nonzero at $t_0$, continuity
would give an interval around $t_0$ on which its absolute value
has a positive lower bound. That interval has positive measure.
Squaring the resulting equality for every real $t$ gives
\eqref{eq:entire-identity} on the critical line.

The function $\Delta_pG^2-F^2$ is entire and vanishes, for example, at
$1/2+i/n$ for every positive integer $n$. These points accumulate
at $1/2$ in its domain. The identity theorem makes the function
zero on $\C$. If $w>0$ almost everywhere, a set of positive
Lebesgue measure has positive $w(t)\,dt$ measure: decompose it
into the sets where $w\geq1/n$. Thus weighted-almost-everywhere
equality implies the equality used in the preceding argument.
\end{proof}

\begin{theorem}[Entire normalization obstruction]
\label{thm:polarization-obstructions:entire}
For $p>1$, there are no nonzero entire functions $F,G$ satisfying
\eqref{eq:boundary-normalization} almost everywhere. In fact no
nonzero entire functions satisfy \eqref{eq:entire-identity}.
\end{theorem}
\begin{proof}
If a nonzero entire function vanishes to infinite order at zero,
its Taylor series is zero in a neighborhood, and the identity
theorem makes it identically zero. Hence nonzero entire $F,G$
have finite orders $m=\ord_0F$ and $n=\ord_0G$ in $\Z_{\geq0}$.
Their local factorizations are
\[
 F(z)=z^m u(z),\qquad G(z)=z^n v(z),\qquad
 u(0)v(0)\ne0.
\]
Lemma~\ref{lem:simple-zero} similarly gives $\Delta_p(z)=zd(z)$ with
$d(0)\ne0$. Equality \eqref{eq:entire-identity} would therefore
force $1+2n=2m$. The left side is odd and the right side is even.
This is impossible. Lemma~\ref{lem:boundary-identity} supplies
\eqref{eq:entire-identity} from the boundary relation.
\end{proof}

The theorem applies to every nonzero element of an entire Sonin
carrier, regardless of its zeros at $0$ or $1$. Adding finite-order
zeros at the Euler lattice cannot remove an odd-order obstruction
to a square. Allowing poles of finite order also does not help:
the local orders then belong to $\Z$, and $1+2n=2m$ remains
impossible. A branch on a cut plane is a different analytic object
and is outside the asserted entire realization.

\begin{corollary}
\label{cor:ambient-polar}
Let $\mathcal L_S=L^2(\R,w_S\,dt)$ and
$\mathcal L_T=L^2(\R,\omega_p w_S\,dt)$, where $w_S>0$ almost everywhere.
The polar isometry of the identity $j:\mathcal L_S\to\mathcal L_T$
is multiplication by $\omega_p^{-1/2}$. It does not restrict to a map
from any nonzero entire-function carrier into itself.
\end{corollary}
\begin{proof}
The bounds on $\omega_p$ make $j$ a bounded bijection. Directly from
the two inner products, $j^*$ is multiplication by $\omega_p$ and
$j^*j=M_{\omega_p}$. Functional calculus gives
$j(j^*j)^{-1/2}=jM_{\omega_p^{-1/2}}$. The boundary output of a
nonzero entire input cannot be the boundary of an entire function
by Theorem~\ref{thm:polarization-obstructions:entire}.
\end{proof}

\subsection{Phases and divisors}

The positive square root in polar decomposition is responsible
for the parity argument. An arbitrary isometry between the ambient
spaces may include a measurable phase $\chi(t)$ of modulus one:
$F_t\mapsto\chi(t)\omega_p(t)^{-1/2}F_t$.
Theorem~\ref{thm:polarization-obstructions:entire} does not exclude
every such phase. It excludes the positive pointwise polar factor
of the specified carrier identity.

The complex reciprocal Euler factor illustrates the difference.
Multiplication by $(1-p^{-z})^{-1}$ has the same boundary modulus
as $\omega_p^{-1/2}$, but its continuation is meromorphic. For an entire
input $F$, the quotient $F/(1-p^{-z})$ is entire exactly when $F$
vanishes at every point $2\pi ik/L$. Necessity follows by comparing
orders at each simple denominator zero. Sufficiency follows by
removing those local singularities; the lattice has no finite
accumulation point.

This cancellation condition is nonempty. For any entire $H$, the
function $F=(1-p^{-z})H$ has the entire quotient $H$. Whether
both belong to a particular Sonin carrier requires that carrier's
growth and norm conditions. The example only establishes the
logical boundary of the positive-square-root obstruction.

There is a local version useful for other candidates. If a
holomorphic function $D$ has odd order at a point $z_0$, then
$DG^2=F^2$ cannot hold for nonzero holomorphic germs $F,G$ at
$z_0$. The preceding proof uses only local factorizations. Thus
one does not need a global formula for the zeros of $D$ once
an odd zero has been found.

For a finite set of distinct primes $E$, the product
$\Delta_E=\prod_{p\in E}\Delta_p$ has an odd zero whenever $E$ is nonempty.
If $|E|$ is odd, zero itself works. If $|E|$ is even, choose
$p\in E$ and $z_0=2\pi i/\log p$. A second prime $q\ne p$
could have a zero there only if $\log q/\log p$ were a positive
integer. This would give $q=p^k$, contrary to distinct primality.
No second-lattice zero has real part zero. The zero of $\Delta_E$
at $z_0$ is therefore simple. The same argument excludes the
positive pointwise normalization for every nonempty finite
prime enlargement, including those with an even number of primes.

\section{Compressed polar decomposition}
\label{sec:compression}

Let $\mathcal L=L^2(\R,w\,dt)$, with $w>0$ almost everywhere,
and let $H\subset\mathcal L$ be closed. Suppose $d$ is real and
$0<c\leq d\leq C<\infty$ almost everywhere. Write
$q_0(x,y)=\langle x,y\rangle_{\mathcal L}$ and
$q_1(x,y)=\langle M_dx,y\rangle_{\mathcal L}$ on $H$.
Equivalent norms give the same closed vector subspace in the
two ambient realizations. Write $H_i=(H,q_i)$ for $i=0,1$.
The identity $J:H_0\to H_1$ is a
bounded bijection between these two Hilbert norms.

\begin{proposition}
\label{prop:polarization-obstructions:compression}
Let $P$ be the $q_0$-orthogonal projection of $\mathcal L$ onto
$H$. The metric operator
\begin{equation}
 A=P M_d|_H
 \label{eq:compressed-metric}
\end{equation}
is positive, self-adjoint and boundedly invertible on $H_0$, with
$c\Id\leq A\leq C\Id$. The restricted identity has
\begin{equation}
 J^*=A,\qquad J^*J=A,\qquad U=JA^{-1/2}.
 \label{eq:intrinsic-polar}
\end{equation}
Here $J^*:H_1\to H_0$ uses the same vector formula $A$ but
different domain and target norms. The $q_1$-orthogonal
projection from the ambient space onto $H$ is
\begin{equation}
 P_1=A^{-1}PM_d.
 \label{eq:changed-projection}
\end{equation}
\end{proposition}
\begin{proof}
For $x,y\in H$, orthogonality of $P$ gives
\[
 q_1(x,y)=q_0(M_dx,y)=q_0(PM_dx,y)=q_0(Ax,y).
\]
Since $M_d$ is self-adjoint, this also equals $q_0(x,Ay)$.
The bounds on $d$ imply
$c\|x\|_0^2\leq q_0(Ax,x)\leq C\|x\|_0^2$.
Thus $A$ is bounded, self-adjoint and bounded below by $c$.
Its range is closed, its kernel is zero, and its range is dense
because its orthogonal complement is $\ker A^*=\ker A$.
It is therefore surjective, with bounded inverse.

The displayed form identity gives the adjoint formula for $J$.
The polar formula follows from the unique positive square root
of $A$. In particular
\[
 q_1(A^{-1/2}x,A^{-1/2}y)
 =q_0(A^{1/2}x,A^{-1/2}y)=q_0(x,y).
\]
Surjectivity of $A^{-1/2}$ makes $U$ a unitary map between
$H_0$ and $H_1$.

For an ambient vector $v$, put $z=A^{-1}PM_dv\in H$.
For every $h\in H$,
\[
 q_1(v-z,h)=q_0(PM_dv-Az,h)=0.
\]
This characterizes $z$ as the $q_1$-orthogonal projection of
$v$. The operator is bounded in the ambient $q_1$ norm by
equivalence of norms, proving \eqref{eq:changed-projection}.
\end{proof}

Applied to $d=\omega_p$ and the closed entire Sonin realization, the
proposition constructs an intrinsic polar isometry. Every output
is an element of the same Hilbert space $H$, so it has the
entire representative belonging to that realization. There is no
conflict with Theorem~\ref{thm:polarization-obstructions:entire}:
the expression $A^{-1/2}$ is not asserted to be $M_{\omega_p^{-1/2}}$.

If $H$ is invariant under $M_d$, its orthogonal complement is
also invariant because $M_d$ is self-adjoint. Then $P$ commutes
with $M_d$, and continuous functional calculus identifies
$A^{-1/2}$ with $M_{d^{-1/2}}|_H$. For $d=\omega_p$, a nonzero
closed subspace consisting of entire boundary functions cannot
have this invariance: its inverse square root would contradict
the entire obstruction. Compression is therefore forced in
any such nonzero realization.

\subsection{A finite-dimensional model}

Even a one-dimensional subspace separates compressed and
pointwise normalization. Take $\mathcal L=\C^2$,
$H=\Span\{(1,1)\}$ and $d=\operatorname{diag}(1,4)$.
For $x=(a,a)$ one has $q_0(x,x)=2|a|^2$ and
$q_1(x,x)=5|a|^2$. The metric operator on $H$ is
$A=(5/2)\Id$. Intrinsic normalization sends
$(a,a)$ to $\sqrt{2/5}(a,a)$. Ambient normalization sends it
to $(a,a/2)$, which is outside $H$ unless $a=0$.
Both operations are isometries on their stated domains.

This example also shows why replacing functional calculus of
a compression by compression of functional calculus is invalid.
On $H$, the operator $PM_{d^{-1/2}}|_H$ is $(3/4)\Id$,
whereas $A^{-1/2}=\sqrt{2/5}\Id$. Their scalar values differ.
No approximation or infinite-dimensional issue is involved.

\subsection{Test maps and symmetry}

An arithmetic test map specifies which geometric vector represents
each compact smooth test in the comparison formula. Changing
that vector changes the pairing being compared with the Weil form.
Suppose $\Gamma_0:\T\to H_0$ and
$\Gamma_1:\T\to H_1$ are complex-linear maps satisfying the
original compatibility $\Gamma_1=J\Gamma_0$. Replacing $J$
by $U$ while retaining both test maps requires
\begin{equation}
 A^{-1/2}\Gamma_0f=\Gamma_0f\qquad(f\in\T).
 \label{eq:test-rigidity}
\end{equation}

\begin{proposition}
\label{prop:test-rigidity}
If $\Gamma_0(\T)$ is dense in $H_0$, the original and polar
test-map compatibilities both hold exactly when $A=\Id$.
Without density they require $A$ to be the identity on the
closure of $\Gamma_0(\T)$.
\end{proposition}
\begin{proof}
The difference $A^{-1/2}-\Id$ is bounded. If it vanishes on
$\Gamma_0(\T)$, it vanishes on its closure. On that closure,
$A^{-1/2}x=x$ implies $A^{1/2}x=x$ by applying $A^{1/2}$,
and then $Ax=x$. Conversely $Ax=x$ implies
$A^{-1/2}x=x$ by functional calculus, or by uniqueness of
the positive square root on the eigenspace for eigenvalue one.
The dense case makes this identity hold on all of $H_0$.
\end{proof}

The proposition is a conditional test. It does not claim that
the compressed Sonin metric operator is nontrivial on any
particular arithmetic test image. A proof of compatibility
must either establish \eqref{eq:test-rigidity}, change the
target test map, or supply a new comparison theorem for the
changed image. Redefining $\Gamma_1=U\Gamma_0$ proves an
analytic identity by construction but does not preserve an
independently defined arithmetic correspondence automatically.

Fourier compatibility can survive compression under verifiable
symmetry hypotheses. Let $R$ be a unitary or antiunitary
operator on $\mathcal L$ that maps $H$ onto $H$ and commutes with
$M_d$. Since $R$ maps $H^\perp$ onto itself, it commutes with
$P$. Therefore $RA=AR$ on $H$. Real polynomial approximation
to $x^{-1/2}$ on $[c,C]$ gives $RA^{-1/2}=A^{-1/2}R$;
real coefficients also cover the antiunitary case.

For example, reflection $v(t)\mapsto v(-t)$ commutes with
$M_{\omega_p}$ because $\omega_p$ is even. It is unitary when the
ambient weight is even. If the carrier is reflection invariant,
the intrinsic normalization respects that symmetry. To call
this the Fourier duality of a semilocal model, one must identify
its actual Fourier operator in the chosen carrier. The shell
transition of Part I has such a published intertwining identity.
The conditional reflection calculation alone does not identify
every geometric or arithmetic duality operator.

\section{Compatibility after compression}
\label{sec:prime-compatibility}

The scalar Euler multipliers for distinct primes commute.
Compression does not preserve this implication for their
metric operators. Let $M_a,M_b$ be commuting self-adjoint
ambient multipliers and $P$ an orthogonal projection. On $H$,
put $A=PM_a|_H$ and $B=PM_b|_H$. Expanding $\Id=P+(\Id-P)$
between the ambient multipliers gives
\begin{equation}
 AB-BA=PM_b(\Id-P)M_a|_H-PM_a(\Id-P)M_b|_H.
 \label{eq:compression-commutator}
\end{equation}
The right side records the two excursions outside $H$.
Commutativity of $M_a,M_b$ does not force either term to vanish.

An explicit example uses the orthonormal vectors
\[
 u=\frac{(1,-1,0)}{\sqrt2},\qquad
 v=\frac{(1,1,-2)}{\sqrt6}
\]
in $\C^3$ and $H=\Span\{u,v\}$. Compress the two coordinate
projections $E_1=\operatorname{diag}(1,0,0)$ and
$E_2=\operatorname{diag}(0,1,0)$. In the basis $(u,v)$ the
compressed matrices are
\begin{equation}
 A_1=\begin{pmatrix}1/2&1/(2\sqrt3)\\1/(2\sqrt3)&1/6\end{pmatrix},
 \qquad
 A_2=\begin{pmatrix}1/2&-1/(2\sqrt3)\\-1/(2\sqrt3)&1/6\end{pmatrix}.
 \label{eq:compressed-matrices}
\end{equation}
Multiplying the matrices yields
\begin{equation}
 A_1A_2-A_2A_1=
 \begin{pmatrix}0&-1/(3\sqrt3)\\1/(3\sqrt3)&0\end{pmatrix}\ne0.
 \label{eq:matrix-commutator}
\end{equation}
Replacing the ambient projections by $\Id+E_1$ and
$\Id+E_2$ makes them positive and invertible without changing
the compressed commutator. These positive diagonal operators
are finite-dimensional models of commuting ambient multiplication
operators; their compressed matrices show the possible failure
of commutation. Thus strict positivity does not repair the failure.

The example is a statement about compression, not a calculation
of the actual Sonin projection. It prevents an argument that
uses scalar commutativity as its only justification for
commuting compressed metric operators. The Sonin operators
may have additional identities; those identities need proof.

The second metric transition must use the inner product
produced by the first. Express both intermediate and final
forms through the base form by setting
$q_p(x,y)=q_0(C_px,y)$ and
$q_{pq}(x,y)=q_0(C_{pq}x,y)$, where
\[
 C_p=PM_{\omega_p}|_H,\qquad
 C_{pq}=PM_{\omega_p\omega_q}|_H.
\]
The metric operator from $q_p$ to $q_{pq}$ is
\begin{equation}
 A_{p,q}=C_p^{-1}C_{pq},
 \label{eq:relative-metric}
\end{equation}
as an operator on the intermediate Hilbert space $(H,q_p)$.
Its positivity and adjoint are determined by $q_p$, as the identity
$q_p(A_{p,q}x,y)=q_0(C_{pq}x,y)=q_{pq}(x,y)$
shows. It need not be self-adjoint for the base form $q_0$.
Its positive square root must therefore be taken in $(H,q_p)$.

Writing every underlying carrier map as the identity, the
polar path through $p$ acts by
$A_{p,q}^{-1/2}C_p^{-1/2}$, and the path through $q$ acts by
$A_{q,p}^{-1/2}C_q^{-1/2}$. A commuting polar square therefore
requires equality of these two operators. The calculation
above neither proves nor disproves that identity for Sonin
spaces. It identifies the operator equation that replaces
commutation of the scalar Euler factors.

There is always an abstract way to choose coherent isometries
relative to one base metric. If $q_S(x,y)=q_0(C_Sx,y)$ on a
common Hilbert carrier, put $V_S=C_S^{-1/2}:H_0\to H_S$ and
$V_{S,T}=V_TV_S^{-1}$. Then
$V_{T,U}V_{S,T}=V_{S,U}$ by cancellation. These maps are
isometries and form commuting squares. They depend on the
chosen base metric and need not be the polar factors of the
one-step carrier identities. Their arithmetic test-map
compatibility is again a separate equation. Abstract coherence
does not construct a local arithmetic correspondence.

\section{Smooth local witnesses}
\label{sec:bumps}

Fix $p>1$, $L=\log p$ and $r=p^{-1/2}$. Choose
$0<\epsilon<L/2$ and a nonzero real function
$\psi\in C_c^\infty(\R)$ supported in $[-\epsilon,\epsilon]$.
For instance one may use
\begin{equation}
 \psi(x)=
 \begin{cases}
 \exp\bigl(-1/(1-(x/\epsilon)^2)\bigr),&|x|<\epsilon,\\
 0,&|x|\geq\epsilon.
 \end{cases}
 \label{eq:explicit-bump}
\end{equation}
All derivatives tend to zero at the endpoints because every
polynomial in $(1-(x/\epsilon)^2)^{-1}$ is dominated by its
negative exponential. Thus this is a compact smooth function,
positive near zero, and $\|\psi\|_2^2>0$.
No unit-mass normalization is required: every local quadratic
value below retains the same positive factor $\|\psi\|_2^2$.

For a finite complex vector $c=(c_0,\ldots,c_n)$, write
\begin{equation}
 f_c(x)=\sum_{j=0}^nc_j\psi(x-jL),\qquad
 C_c(z)=\sum_{j=0}^nc_jz^j.
 \label{eq:bump-packet}
\end{equation}
The translates have disjoint supports. Hence
\begin{equation}
 \|f_c\|_2^2=\|\psi\|_2^2\sum_{j=0}^n|c_j|^2.
 \label{eq:packet-norm}
\end{equation}
In particular a nonzero coefficient vector produces a nonzero
test, without any possible cancellation between adjacent bumps.

\begin{proposition}
\label{prop:polarization-obstructions:bumps}
For the packet \eqref{eq:bump-packet},
\begin{align}
 F_{f_c}(z)&=F_\psi(z)C_c(e^{Lz}),
 \label{eq:packet-moments}\\
 W_p(f_c*f_c^*)&=L\|\psi\|_2^2Q_r(c),
 \label{eq:packet-local}\\
 Q_r(c)&=\sum_{i\ne j}c_i\overline{c_j}r^{|i-j|}.
 \label{eq:toeplitz}
\end{align}
Consequently $C_c(r)=C_c(r^{-1})=0$ implies $f_c\in\T_0$.
\end{proposition}
\begin{proof}
Changing variables $u=x-jL$ in the integral for the Laplace
transform gives
$F_{\psi(\cdot-jL)}(z)=e^{jLz}F_\psi(z)$.
Summing the finitely many terms proves
\eqref{eq:packet-moments} for every complex $z$.
At $z=1/2$ and $z=-1/2$ the polynomial arguments are
$r^{-1}$ and $r$, respectively.

Put $h_\psi=\psi*\psi^*$. Its support lies in
$[-2\epsilon,2\epsilon]$ and $h_\psi(0)=\|\psi\|_2^2$.
Expanding the finite convolution sum gives
\[
 (f_c*f_c^*)(x)=\sum_{i,j=0}^nc_i\overline{c_j}
 h_\psi(x-(i-j)L).
\]
At $x=kL$ with integer $k$, every term vanishes unless
$k=i-j$, because $2\epsilon<L$. Therefore, for $k\geq1$,
\begin{align*}
 (f_c*f_c^*)(kL)&=\|\psi\|_2^2\sum_{i-j=k}c_i\overline{c_j},\\
 (f_c*f_c^*)(-kL)&=\|\psi\|_2^2\sum_{j-i=k}c_i\overline{c_j}.
\end{align*}
These sums are zero for $k>n$. Substitution in the actual
local distribution \eqref{eq:local-Weil} now gives
\eqref{eq:packet-local}. Each ordered pair $i\ne j$ occurs
once, with its weight $r^{|i-j|}$. This also shows directly
that $Q_r(c)$ is real, since reversing $i,j$ conjugates a term.
\end{proof}

The condition on the support of $\psi$ is exact. There is no
limit as $\epsilon$ tends to zero, and the normalization
$\|\psi\|_2^2$ is strictly positive. All prime powers that
can contribute to the packet are included in
\eqref{eq:packet-local}. Terms at larger powers vanish by
support, rather than being discarded as small.

Let $a_r=r+r^{-1}$ and consider
\begin{equation}
 P_r(z)=1-a_rz+z^2=(z-r)(z-r^{-1}).
 \label{eq:moment-factor}
\end{equation}
Any coefficient polynomial divisible by $P_r$ has both
required roots. This provides a direct way to impose the pole
conditions before testing the sign, without projecting a
numerical eigenvector onto an approximate moment kernel.

\begin{lemma}
\label{lem:short-values}
For $0<r<1$ and $a_r=r+r^{-1}$,
\begin{align}
 Q_r(1,-a_r,1)&=-4-2r^2,
 \label{eq:negative-value}\\
 Q_r(1,1-a_r,1-a_r,1)&=-8-4r^2+6r+2/r.
 \label{eq:positive-value}
\end{align}
The respective coefficient polynomials are $P_r(z)$ and
$P_r(z)(1+z)$.
\end{lemma}
\begin{proof}
For real coefficients one can sum by distance:
\[
 Q_r(c)=2\sum_{k=1}^nr^k\sum_{j=0}^{n-k}c_{j+k}c_j.
\]
For the three-term vector the inner sums are $-2a_r$ at
distance one and $1$ at distance two. Thus
$Q=-4a_rr+2r^2=-4-2r^2$.

For the four-term vector set $b=1-a_r$. The inner sums are
$2b+b^2$ at distance one, $2b$ at distance two, and $1$ at
distance three. Hence
\begin{align*}
 Q&=2r(2b+b^2)+4r^2b+2r^3\\
  &=2r(3-4a_r+a_r^2)+4r^2(1-a_r)+2r^3.
\end{align*}
Using $a_rr=1+r^2$ and
$a_r^2=r^2+2+r^{-2}$ reduces this expression to
\eqref{eq:positive-value}. Multiplication of $P_r$ by
$1+z$ gives coefficients $(1,1-a_r,1-a_r,1)$, which proves
the polynomial statements.
\end{proof}

\begin{theorem}[Exact signs at eleven]
\label{thm:polarization-obstructions:signs}
There exist nonzero real tests $f_-,f_+\in\T_0$ such that
\begin{align}
 W_{11}(f_-*f_-^*)&=-\frac{46}{11}(\log11)\|\psi\|_2^2<0,
 \label{eq:eleven-minus}\\
 W_{11}(f_+*f_+^*)&=\frac{28\sqrt{11}-92}{11}
 (\log11)\|\psi\|_2^2>0.
 \label{eq:eleven-plus}
\end{align}
They can be chosen as three and four translates of any bump
\eqref{eq:explicit-bump} with $2\epsilon<\log11$.
\end{theorem}
\begin{proof}
Set $r=1/\sqrt{11}$, so $a_r=12/\sqrt{11}$, and use the
two coefficient vectors in Lemma~\ref{lem:short-values}.
Their polynomials are divisible by $P_r$, so
Proposition~\ref{prop:polarization-obstructions:bumps} gives
both pole moments exactly zero. Their first coefficient is
one, and \eqref{eq:packet-norm} makes both tests nonzero.

Substituting $r^2=1/11$ into
\eqref{eq:negative-value} gives $-46/11$. Substituting
$r=1/\sqrt{11}$ into \eqref{eq:positive-value} gives
$(28\sqrt{11}-92)/11$. This is positive because
$7\sqrt{11}>23$: both sides are positive and their squares
satisfy $49\cdot11=539>529=23^2$.
Finally $\log11>0$ and $\|\psi\|_2^2>0$, so the signs
are unchanged in the smooth distribution values.
\end{proof}

Translation of the complete packet by a common real number
does not change its autocorrelation. It multiplies each
Laplace moment by a nonzero exponential and therefore
preserves a zero moment. The negative packet can consequently
be centered in $[-L-\epsilon,L+\epsilon]$, and the positive
packet in $[-3L/2-\epsilon,3L/2+\epsilon]$. This centering
will make the support condition in the geometric obstruction
explicit.

\section{Finite packets at every prime}
\label{sec:all-primes}

The short positive packet at $11$ is sufficient for the
orthogonal-extension obstruction. A longer packet gives the
same two-sign statement at every prime. Its exact formula
also explains why moment cancellation does not restore the
sign of a local distribution.

\begin{proposition}
\label{prop:packet-formula}
For $0<r<1$, $N\geq2$, let $c^{(N)}$ be the coefficient
vector of
\begin{equation}
 C_N(z)=P_r(z)(1+z+\cdots+z^{N-1}).
 \label{eq:long-packet}
\end{equation}
Then
\begin{equation}
 Q_r(c^{(N)})=\frac{2N(1-r)^3}{r}+4-6r-\frac2r.
 \label{eq:long-value}
\end{equation}
\end{proposition}
\begin{proof}
On the unit circle $z=e^{i\theta}$ the absolutely convergent
Fourier series of the off-diagonal Toeplitz kernel is
\begin{align}
 k_r(\theta)
 &=\sum_{k\ne0}r^{|k|}e^{ik\theta}
 =\frac{1-r^2}{1-r(z+z^{-1})+r^2}-1.
 \label{eq:toeplitz-symbol}
\end{align}
Expanding a finite polynomial $C$ and integrating its product
with this series term by term gives
\begin{equation}
 Q_r(c)=\frac1{2\pi}\int_{-\pi}^{\pi}
 k_r(\theta)|C(e^{i\theta})|^2\,d\theta.
 \label{eq:symbol-form}
\end{equation}
Absolute convergence justifies the interchange, and
orthogonality of integer exponentials leaves exactly the
coefficient sum in \eqref{eq:toeplitz}.

Put $s=z+z^{-1}$ and $D=1-rs+r^2$. The moment polynomial
in \eqref{eq:moment-factor} satisfies
\[
 P_r(z)=z(s-r-r^{-1})=-zD/r.
\]
For $|z|=1$, $D$ is real, so $|P_r(z)|^2=D^2/r^2$.
Multiplying \eqref{eq:toeplitz-symbol} by this square cancels
the denominator:
\begin{align}
 k_r(\theta)|P_r(z)|^2
 &=\frac{(rs-2r^2)D}{r^2}\\
 &=(-4-2r^2)+(3r+r^{-1})(z+z^{-1})-(z^2+z^{-2}).
 \label{eq:laurent-cancellation}
\end{align}
For completeness, expansion of the first line yields
$(r^{-1}+3r)s-s^2-2-2r^2$, and
$s^2=z^2+2+z^{-2}$ gives the second line.

The polynomial $B_N(z)=1+\cdots+z^{N-1}$ satisfies
\[
 |B_N(z)|^2=\sum_{k=-(N-1)}^{N-1}(N-|k|)z^k.
\]
The constant coefficient of
$k_r|P_r|^2|B_N|^2$ is therefore
\[
 N(-4-2r^2)+2(N-1)(3r+r^{-1})-2(N-2).
\]
This formula includes $N=2$, when the final contribution is
zero. Simplifying its coefficient of $N$ gives
$2(r^{-1}-3+3r-r^2)=2(1-r)^3/r$, and its constant term
is $4-6r-2/r$. Equation~\eqref{eq:symbol-form} proves
\eqref{eq:long-value}.
\end{proof}

\begin{corollary}
\label{cor:all-prime-signs}
For every prime $p$, the quadratic form $f\mapsto W_p(f*f^*)$
has both positive and negative values on nonzero real
tests in $\T_0$.
\end{corollary}
\begin{proof}
The three-term packet always gives $-4-2r^2<0$.
The coefficient of $N$ in \eqref{eq:long-value} is strictly
positive for $0<r<1$. Choose an integer $N\geq2$ such that
\[
 N>\frac{r(6r+2/r-4)}{2(1-r)^3}.
\]
Then \eqref{eq:long-value} is positive. The polynomial
$C_N$ is divisible by $P_r$ and has first coefficient one,
so its smooth packet is nonzero and belongs to $\T_0$.
Proposition~\ref{prop:polarization-obstructions:bumps}
transfers the exact positive coefficient value to $W_p$.
\end{proof}

The support of the longer positive packet grows with $N$.
The corollary therefore gives a local sign result on the union
of compact test supports. It gives no common small support
on which every prime is visible. On a fixed support, all
sufficiently large primes have identically zero local form.
This restriction is part of the arithmetic geometry of the
test system, rather than a defect in the packet construction.

The Toeplitz symbol itself is positive at $\theta=0$ and
negative at $\theta=\pi$. That observation would suggest
frequency localization, but it would not by itself impose
the two moment conditions. The finite factor $P_r$ imposes
those conditions exactly. Its roots $r,r^{-1}$ are off the
unit circle, so the factor leaves nontrivial Fourier packets
available. The cancellation in \eqref{eq:laurent-cancellation}
turns that observation into the finite proof above.

\section{The constrained finite form}
\label{sec:finite-form}

The four-bump example at $11$ has a useful minimality property
within this specific translated-bump family. The two moment
conditions reduce the coefficient space to a polynomial
divisibility condition, and the restricted quadratic form can
be diagonalized exactly. This is a finite-dimensional statement
about equally spaced translates of the same bump. It is not
a minimal-support theorem for arbitrary compact smooth tests.

\begin{proposition}
\label{prop:four-bump-signature}
Let $0<r<1$, and let $c=(c_0,c_1,c_2,c_3)\in\C^4$ satisfy
$C_c(r)=C_c(r^{-1})=0$. There are unique
$\alpha,\beta\in\C$ such that
\begin{equation}
 C_c(z)=P_r(z)(\alpha+\beta z).
 \label{eq:four-factorization}
\end{equation}
On these coordinates the constrained form is
\begin{equation}
 Q_r(c)=b_0(|\alpha|^2+|\beta|^2)
           +2b_1\operatorname{Re}(\alpha\overline\beta),
 \quad b_0=-4-2r^2,\quad b_1=3r+r^{-1}.
 \label{eq:restricted-matrix}
\end{equation}
Its eigenvalues in the coefficient coordinates
$(\alpha,\beta)$ are $b_0+b_1$ and $b_0-b_1$.
\end{proposition}
\begin{proof}
The roots $r$ and $r^{-1}$ are distinct, so their two
linear factors divide $C_c$. Since $C_c$ has degree at
most three and $P_r$ is monic of degree two, the quotient
has degree at most one and is unique. This proves
\eqref{eq:four-factorization}, including the zero polynomial.
Substitute $C_c=P_r(\alpha+\beta z)$ into
\eqref{eq:symbol-form} and use
\eqref{eq:laurent-cancellation}. The product
\[
 |\alpha+\beta z|^2=|\alpha|^2+|\beta|^2
 +\overline\alpha\beta z+\alpha\overline\beta z^{-1}
\]
has frequencies only zero and one. Its constant coefficient
after multiplication by the Laurent polynomial is the
right side of \eqref{eq:restricted-matrix}; the terms at
frequency two contribute nothing. The matrix is
$\left(\begin{smallmatrix}b_0&b_1\\b_1&b_0\end{smallmatrix}\right)$.
Its eigenvectors $(1,1)$ and $(1,-1)$ have the stated
eigenvalues.
\end{proof}

The word eigenvalue here refers to the Euclidean norm of
the quotient coefficients $(\alpha,\beta)$. The original
four coefficients have a different induced norm under
\eqref{eq:four-factorization}. The signature, which is what
the sign test uses, is unchanged by this injective change
of coordinates.

\begin{corollary}
\label{cor:smallest-prime}
Among arithmetic primes, $11$ is the smallest prime for
which four equally spaced bumps can give a positive local
value while their coefficient polynomial has both moment
roots. For $p\leq7$ the form on that two-dimensional
coefficient kernel is negative definite. For every prime
$p\geq11$ its signature is $(1,1)$.
\end{corollary}
\begin{proof}
The eigenvalue $b_0-b_1$ is negative for every $0<r<1$.
The sign of $b_0+b_1$ is the sign of
\[
 g(r)=r(b_0+b_1)=1-4r+3r^2-2r^3.
\]
Its derivative is
$g'(r)=-4+6r-6r^2=-6(r-1/2)^2-5/2<0$.
At $r=1/\sqrt{11}$, the value of $b_0+b_1$ is
$(14\sqrt{11}-46)/11>0$, by the integer-square
comparison used in Theorem~\ref{thm:polarization-obstructions:signs}.
It remains positive for smaller $r$, which includes every
prime $p\geq11$.

At $r=1/\sqrt7$,
\[
 g(r)=\frac{10\sqrt7-30}{7\sqrt7}<0,
\]
because $\sqrt7<3$. It remains negative for larger $r$,
which includes every prime at most seven. The only primes
below eleven are $2,3,5,7$; the intervening integers
$8,9,10$ are composite. This proves all the assertions.
\end{proof}

Three equally spaced bumps cannot give a positive value
under the same coefficient moment conditions, at any prime.
Their polynomial has degree at most two, so it must be
$\alpha P_r$. Its quadratic value is
$|\alpha|^2(-4-2r^2)$, strictly negative unless the test is
zero. With two bumps or one bump, a polynomial of degree
at most one with both distinct roots is zero. Thus the
three-term negative witness and the four-term positive
witness at $11$ use the smallest possible numbers of
translates within this coefficient construction.

For the explicit nonnegative bump
\eqref{eq:explicit-bump}, both values $F_\psi(1/2)$ and
$F_\psi(-1/2)$ are strictly positive. Equation
\eqref{eq:packet-moments} then makes the coefficient
root conditions equivalent to the actual two zero moments.
The minimality statement applies to all packets built from
this bump, rather than just to a sufficient algebraic
condition. If a different bump already has zero moments,
the equivalence can fail, and this argument does not give
the same minimality statement for that family.

\section{Orthogonal primitive extensions}
\label{sec:extensions}

Let $\mathcal D$ be a complex test space on which $B_p$ is
defined. It may be $\T_0$ or a subspace containing the
specific witness under consideration. Let $P_S,P_T$ be
complex vector spaces equipped with Hermitian forms
$I_S,I_T$. No completeness assumption is needed in the
following algebraic argument.

\begin{definition}
\label{def:polarization-obstructions:extension}
An orthogonal nonpositive extension on $\mathcal D$ consists
of an isometric complex-linear inclusion
$\iota:P_S\to P_T$ and complex-linear maps
$\Gamma_S:\mathcal D\to P_S$,
$\Gamma_T,\eta:\mathcal D\to P_T$ such that
\begin{align}
 \Gamma_Tf&=\iota\Gamma_Sf+\eta f,
 \label{eq:extension-decomp}\\
 I_T(\iota x,\eta f)&=0\qquad(x\in P_S,f\in\mathcal D),
 \label{eq:extension-orthogonal}\\
 I_T(\eta f,\eta f)&\leq0\qquad(f\in\mathcal D).
 \label{eq:extension-negative}
\end{align}
The isometry condition means
$I_T(\iota x,\iota y)=I_S(x,y)$.
\end{definition}

Nonpositivity of all of $I_T$ implies
\eqref{eq:extension-negative}, but the definition only needs
the sign on the added test image. Thus the exclusion below
also applies when the ambient geometric pairing has other
signs outside the relevant sector.

\begin{theorem}[Orthogonal extension obstruction]
\label{thm:polarization-obstructions:extension}
Suppose $\mathcal D$ contains the positive test $f_+$ in
Theorem~\ref{thm:polarization-obstructions:signs}. There is
no orthogonal nonpositive extension on $\mathcal D$ for
which the one-prime comparison at $p=11$ is
\begin{equation}
 I_T(\Gamma_Tf,\Gamma_Tg)-I_S(\Gamma_Sf,\Gamma_Sg)
 =W_{11}(f*g^*)\qquad(f,g\in\mathcal D).
 \label{eq:target-increment}
\end{equation}
For an arbitrary prime $p$, the same conclusion holds on
$\mathcal D=\T_0$ with $W_p$ in place of $W_{11}$.
\end{theorem}
\begin{proof}
Expand \eqref{eq:extension-decomp} in the Hermitian form.
The old-old term cancels by isometry. The two mixed terms
are zero by \eqref{eq:extension-orthogonal} and Hermitian
symmetry. Therefore
\[
 I_T(\Gamma_Tf,\Gamma_Tf)-I_S(\Gamma_Sf,\Gamma_Sf)
 =I_T(\eta f,\eta f)\leq0.
\]
At $f=f_+$ this contradicts the strict positivity in
\eqref{eq:eleven-plus}. For every prime, use the positive
test supplied by Corollary~\ref{cor:all-prime-signs}.
\end{proof}

The proof tests a prescribed equation and sign together.
It does not assume an intersection form exists and then
derive a contradiction from the existence alone. Removing
orthogonality, allowing corrections in the comparison, or
requiring sign only at a later support-complete stage changes
the hypotheses of the theorem.

The opposite local sign is useful as well. If a proposed
orthogonal sector were nonnegative and its prescribed
increment were $+W_p$, the negative packet would exclude
that version. If the intended comparison used $-W_p$, the
roles of the two packets would interchange. The sign issue
is not repaired by choosing one uniform sign for an
orthogonal local summand on all tests in $\T_0$.

\subsection{Mixed pairings}

\begin{proposition}
\label{prop:polarization-obstructions:mixed}
Suppose $\iota:P_S\to P_T$ is isometric and
$\Gamma_T=\iota\Gamma_S+\eta$, without orthogonality.
Then the diagonal increment is
\begin{equation}
 I_T(\Gamma_Tf,\Gamma_Tf)-I_S(\Gamma_Sf,\Gamma_Sf)
 =2\operatorname{Re}I_T(\iota\Gamma_Sf,\eta f)
 +I_T(\eta f,\eta f).
 \label{eq:mixed-defect}
\end{equation}
If isometry is also dropped, one must add
\begin{equation}
 E(f,g)=I_T(\iota\Gamma_Sf,\iota\Gamma_Sg)
       -I_S(\Gamma_Sf,\Gamma_Sg)
 \label{eq:old-defect}
\end{equation}
to the corresponding sesquilinear expansion.
\end{proposition}
\begin{proof}
Sesquilinearity expands the pairing of
$\iota\Gamma_Sf+\eta f$ with itself into four terms.
The two mixed terms are complex conjugates, giving the real
part in \eqref{eq:mixed-defect}. The old-old term cancels
under isometry. Without isometry its remainder is exactly
\eqref{eq:old-defect}. The same expansion for two different
tests gives the sesquilinear statement.
\end{proof}

Even a negative definite one-dimensional form permits a
positive increment through mixed terms. On $\C$, take
$I(x,y)=-x\overline y$, old vector $x=1$ and added vector
$y=-1/2$. Then
\[
 I(x+y,x+y)-I(x,x)=-1/4+1=3/4>0,
\]
although $I(y,y)=-1/4<0$. Formula
\eqref{eq:mixed-defect} gives a mixed contribution of $1$.
This explicit example rules out any inference that a
nonpositive total pairing makes every transition increment
nonpositive.

There is a quantitative constraint if the whole form is
nonpositive. Applying Cauchy-Schwarz to the positive
semidefinite form $-I_T$ gives
\[
 |I_T(x,y)|^2\leq(-I_T(x,x))(-I_T(y,y)).
\]
For completeness, this inequality follows by requiring
$-I_T(x+zy,x+zy)\geq0$ for all complex $z$ and minimizing
the resulting quadratic polynomial; if $I_T(y,y)=0$, the
linear term must vanish. Thus mixed terms must obey a bound,
but they can have either sign and can dominate the added
diagonal term. A geometric construction must calculate them
before applying any sign argument.

\section{Support-indexed comparison}
\label{sec:support}

Part I, Definition 12.1,
uses pairs $(S,a)$ with all primes at most $e^{2a}$ included
in $S$. Let $\T_a$ denote its tests supported in $[-a,a]$.
The following locality statement explains exactly when the
orthogonal extension obstruction can be applied.

\begin{proposition}
\label{prop:polarization-obstructions:locality}
If $f,g\in\T_a$ and $p>e^{2a}$, then $W_p(f*g^*)=0$.
Consequently, if $(S,a)$ is support admissible and $p\notin S$,
adjoining $p$ has zero local Weil increment on $\T_a$.
\end{proposition}
\begin{proof}
If $(f*g^*)(x)\ne0$, there are points $u$ and $v$ in the
supports of $f$ and $g$ with $x=u-v$. Thus
$\supp(f*g^*)\subset[-2a,2a]$. If $p>e^{2a}$, then
$k\log p>2a$ for every $k\geq1$. Both evaluations in
\eqref{eq:local-Weil} vanish. At a support-admissible pair,
every prime outside $S$ satisfies this strict inequality.
\end{proof}

For the centered positive witness at $11$, one may take
$a=3\log11/2+\epsilon$. A support-complete place set
therefore contains all primes at most $11^3e^{2\epsilon}$,
including $11$. The transition that adds $11$ while testing
this function starts from an incomplete set for this support.
Theorem~\ref{thm:polarization-obstructions:extension} concerns
the stronger demand that this intermediate transition itself
have an orthogonal nonpositive added image and exact increment.

In a support-indexed construction, an admissible enlargement
from $(S,a)$ to $(T,b)$ with $b>a$ admits new tests. A prime
already irrelevant at radius $a$ may contribute at radius $b$.
The comparison on the old tests and the construction on the
new tests are different obligations. Fixing a test and
telescoping local distributions across the primes in $T\setminus S$
is valid algebraically, but intermediate geometric sectors
need not satisfy the final sign theorem.

At the level of the local distributions, path independence
is exact. For finite disjoint sets of primes $E,F$,
\[
 \sum_{p\in E\cup F}W_p(h)=\sum_{p\in E}W_p(h)
 +\sum_{p\in F}W_p(h).
\]
This identity follows from a finite sum. If geometric
transition defects are to reproduce it, their mixed and
boundary terms must telescope to the same result. It does
not follow that each intermediate diagonal defect has a
fixed sign.

\subsection{A finite support bound}

There is a simple continuity estimate for each local form.
For $f,g\in\T_a$, Cauchy-Schwarz applied to their
correlation gives
$|(f*g^*)(x)|\leq\|f\|_2\|g\|_2$ for every $x$.
Let $M=\lfloor2a/L\rfloor$. Then
\begin{equation}
 |W_p(f*g^*)|
 \leq2L\left(\sum_{k=1}^{M}r^k\right)\|f\|_2\|g\|_2
 \leq\frac{2Lr}{1-r}\|f\|_2\|g\|_2.
 \label{eq:local-bound}
\end{equation}
The empty sum is zero. The estimate extends the local form
continuously to $L^2$ tests supported in the same interval.
It also proves continuity on every compact test stage.

Only finitely many primes contribute on $\T_a$, so the
sum of their local forms is continuous on that stage.
This does not provide a uniform bound as $a$ grows, nor
does it address the singular archimedean subtraction.
Those two issues belong to completion of the full comparison.
The local sign witnesses and their obstruction require
neither an infinite Euler product nor a global norm limit.

\section{Pole moments and completion}
\label{sec:completion}

Removing the two pole moments is a useful way to test a
primitive ansatz, but it does not give a dense subspace
from which the full sign follows by continuity. The moment
functionals are continuous in the compact smooth topology.
They are also continuous in $L^2([-a,a])$, since
\[
 |F_f(\pm1/2)|
 \leq\|f\|_2\left(\int_{-a}^ae^{\pm x}\,dx\right)^{1/2}.
\]
Their common kernel is therefore closed in either setting.

\begin{proposition}
\label{prop:moment-complement}
On $\T$, the two functionals $f\mapsto F_f(1/2)$ and
$f\mapsto F_f(-1/2)$ are linearly independent and have
a two-dimensional smooth complement. For any $a>0$ the
same statement holds on $\T_a$.
\end{proposition}
\begin{proof}
Choose a nonnegative nonzero compact smooth bump $\beta$
whose support and a nonzero translate by $b$ are both
contained in $(-a,a)$. This is possible by first taking
the support sufficiently small. Let $\beta_0=\beta$ and
$\beta_1(x)=\beta(x-b)$. Set
$u=F_\beta(1/2)>0$ and $v=F_\beta(-1/2)>0$.
The matrix of their two moments is
\begin{equation}
 M=\begin{pmatrix}u&e^{b/2}u\\v&e^{-b/2}v\end{pmatrix}.
 \label{eq:moment-matrix}
\end{equation}
Its determinant is $uv(e^{-b/2}-e^{b/2})\ne0$ because
$b\ne0$. Thus the moment map is onto $\C^2$ on
$\Span\{\beta_0,\beta_1\}$. For each $f$, solve
\[
 M\begin{pmatrix}c_0\\c_1\end{pmatrix}
 =\begin{pmatrix}F_f(1/2)\\F_f(-1/2)\end{pmatrix}.
\]
Then $f-c_0\beta_0-c_1\beta_1$ has both moments zero,
and the decomposition is unique because $M$ is invertible.
All functions remain in $\T_a$. The same construction
works in $\T$ without a fixed support bound.
\end{proof}

Suppose a Hermitian form $B$ has a known sign on $\T_0$.
The decomposition from Proposition~\ref{prop:moment-complement}
gives a block form on
$\T_0\oplus\Span\{\beta_0,\beta_1\}$. A sign on the
first block leaves the finite degree block and its mixed
entries undetermined. For $f=f_0+v$ the exact identity is
\[
 B(f,f)=B(f_0,f_0)+2\operatorname{Re}B(f_0,v)+B(v,v).
\]
This is why a comparison proved only after killing moments
needs a separate treatment of degree and pole corrections
before the all-test Weil criterion can be invoked.

\subsection{Limits of finite norm comparisons}

For a finite prime set $E$, the local norm bounds multiply:
\begin{equation}
 \prod_{p\in E}(1-p^{-1/2})^2q_S(F,F)
 \leq q_{S\cup E}(F,F)
 \leq\prod_{p\in E}(1+p^{-1/2})^2q_S(F,F).
 \label{eq:finite-bounds}
\end{equation}
These are finite inequalities. Without a separate estimate
they do not give constants uniform in $E$. An arithmetic
completion might use another topology, renormalized pairings
or a support-indexed limit. The needed estimate must be
proved in that stated topology.

A related operator distinction can be established abstractly.
Let $H$ be a Hilbert space, $P_n$ proper orthogonal
projections, and $T_n=P_nA_nP_n$ bounded operators on $H$.
If $T$ is bounded below by $c>0$, then
\begin{equation}
 \|T-T_n\|\geq c\qquad\hbox{for every }n.
 \label{eq:norm-limit-obstruction}
\end{equation}
Indeed choose a unit vector $v_n\in\ker P_n$.
Then $T_nv_n=0$, while $\|Tv_n\|\geq c$.
Thus zero-extended proper cutoffs cannot converge in
operator norm to a bounded-below operator.

Strong convergence can still hold. Suppose $T_n$ are
uniformly bounded by $M$ and $T_nx$ converges for every
$x$ in a dense linear subspace $D\subset H$.
For arbitrary $x\in H$, choose $y\in D$ with
$\|x-y\|<\epsilon/(4M)$ when $M>0$.
For sufficiently large $n,m$,
\[
 \|T_nx-T_mx\|
 \leq2M\|x-y\|+\|T_ny-T_my\|<\epsilon.
\]
Completeness gives $Tx=\lim_nT_nx$; the limit is linear
and bounded by $M$. The case $M=0$ is trivial.
This argument specifies sufficient hypotheses for a
strong completion, but neither hypothesis follows merely
from the one-prime bounds in \eqref{eq:finite-bounds}.

\subsection{Dynamics and domains}

The usual Sonin condition imposes simultaneous spatial
and Fourier gaps. With the Fourier convention
$\mathcal Ff(\xi)=\int f(x)e^{-2\pi ix\xi}\,dx$, the
unitary dilation $D_bf(x)=b^{1/2}f(bx)$ satisfies
\[
 \mathcal F(D_bf)(\xi)=b^{-1/2}\mathcal Ff(\xi/b).
\]
If both original gaps have radius $\lambda$, the transformed
spatial gap has radius $\lambda/b$ and the Fourier gap
has radius $b\lambda$. These formulae follow by substitution
in the support conditions and in the Fourier integral.
They do not establish invariance of a fixed Sonin space
under the full dilation group.

A weight-one adjoint identity also has a domain condition.
If a densely defined operator $\Theta$ satisfies
$\Dom\Theta^*=\Dom\Theta$ and
$\Theta^*=\Id-\Theta$ on that domain, then
$A=(\Theta-\frac12\Id)/i$ is self-adjoint.
To check this, set $B=\Theta-\frac12\Id$. The bounded
shift leaves the adjoint domain unchanged and gives
$B^*=-B$. Since $A=-iB$, the scalar adjoint rule yields
$A^*=iB^*=-iB=A$ on the same domain.
An equality of differential expressions on a test core
does not imply this domain equality.

This conditional operator statement does not supply a
realization of the zeta zeros. A spectral argument would
also have to identify every nontrivial zero, with its
multiplicity, in the stated operator realization. The
entire normalization and local sign arguments above do
not make that identification.

\section{Arithmetic geometry and prior work}
\label{sec:prior}

The geometric motivation has several distinct mathematical
sources. Semilocal adelic quotients and their crossed-product
algebras provide arithmetic actions and local trace formulae.
Connes's semilocal cutoff theorem retains a divergent term
and Fourier-normalized principal values
\cite[Section VII, Theorem 4]{C98}. His global trace condition
in the following section has an RH-equivalent status.
The finite trace theorem therefore cannot be used as an
unconditional global positivity theorem.

Connes, Consani and Marcolli construct a descended trace
pairing through a restriction-map cokernel and compare it
with primitive cohomology
\cite[Sections 4.2 to 5]{CCM07}. Their realization includes
the full zero divisor with multiplicities, and their
positivity criterion separates the trace identity from
the sign \cite[Theorem 6.1 and Proposition 6.2]{CCM07}.
This provides a closer arithmetic precedent than an
abstract Hilbert-space normalization. The present extension
obstruction concerns an extra orthogonality requirement;
it is not a contradiction of that trace construction.

The semilocal Sonin isomorphisms in \cite[Theorem 4.6]{CCM24}
are topological isomorphisms, with distinct inherited
Hilbert norms in the entire realization. Earlier work of
Burnol connects Sonin spaces, co-Poisson summation and
completed Mellin transforms
\cite[Introduction and Section 2]{B04}.
The common entire carrier is therefore prior analytic
structure. The elementary contribution here is the
odd-order test for the particular positive multiplier
proposed as a normalization of its metric transition.

A sheaf organization supplies another kind of compatibility.
Connes's account places semilocal crossed products over
$\operatorname{Spec}\Z$ and describes their cohomological
setting \cite[Section 7.3, Theorem 7.2]{C26}.
Restriction maps and a generic stalk specify how local
arithmetic data fit together. A Hermitian pairing requires
further duality and trace data, followed by a sign theorem.
An identification of the resulting complex with the analytic
Sonin carrier must respect those data and the trace comparison.

The adelic divisor construction of Connes and Consani uses
rank-one torsion-free groups with possibly degenerate
seminorms and identifies their classes as a Picard monoid
\cite[Definition 3.1 and Theorem 3.4]{CC26}.
The framed uniformization in that work provides geometric
coordinates, not an arithmetic Hodge-index theorem for the
Weil test form. In particular a monoid allowing singular
metrics cannot be treated as an already polarized abelian
variety without an additional construction and proof.

Deninger's proposed arithmetic cohomology makes the desired
sign and adjoint relation explicit through a Hodge-star
pairing \cite[Section 3]{D05}. A related weight identity is
proved for a foliated three-manifold under stated geometric
hypotheses \cite[Theorem 4.1]{D05}. The arithmetic realization
is a further conjectural step. This comparison explains why
an independently constructed pairing is essential: once the
positive pairing and the complete spectral realization are
assumed, the centered spectral conclusion is a short
operator argument.

These sources motivate a construction through arithmetic
duality rather than prescribing the two unsuccessful
repairs studied here. The local sign theorem excludes
orthogonal signed increments under its stated hypotheses.
It leaves open a geometric pairing whose primitive
projection changes with support, whose intermediate forms
are indefinite, or whose mixed terms reproduce local
arithmetic contributions.

\subsection{Chain maps and Fourier duality}

Part I, Proposition 14.1, supplies a concrete algebraic
repair of the nonmultiplicative shell tensor map. The two
shell indicators $e_0,e_1$ individually induce algebra
embeddings. Their signed combination extends to the
Hochschild complex by taking the same linear combination
of the two induced chain maps in every degree. This is
different from tensoring the signed degree-zero map with
itself. Its chain compatibility follows from
multiplicativity of the two individual embeddings.

The corrected map also has a finite-prime chain square,
as established in Part I, Section 15. These facts supply
an algebraic correspondence before a primitive projection
or pairing has been constructed. The maps act on
uncompleted algebraic Hochschild chains. Applying a
topological trace realization requires a specified
completion and continuity statements for these maps.

The repair does not preserve the original componentwise
Fourier symmetry in positive chain degree. Part I,
Proposition 14.2, evaluates its degree-one coefficient
on the small ball $B_{-1}^2$. The coefficient itself
vanishes there, whereas its local Fourier transform is
$-(p-1)^3/p^2$. In degree $n\geq1$, Part I, equation (92),
gives the discrepancy
\[
 (p-1)^{n+1}\bigl(p^{-(n+1)}-p^{-1}\bigr)<0
\]
on $B_{-1}^{n+1}$. Both values are exact local-field
calculations. The higher-degree failure cannot be repaired
by citing the Fourier-fixed degree-zero shell alone.

This comparison distinguishes two uses of the word
duality. On a Hilbert carrier, a bounded symmetry may
commute with a compressed metric operator under the
hypotheses in Section~\ref{sec:compression}. On a
Hochschild complex, the relevant duality must interact
with its algebra product and boundary. Local Fourier
transformation exchanges pointwise multiplication and
additive convolution; it is not being assumed to be a
chain map for either unmodified complex. A correspondence
that is correct for one product therefore needs an
explicit comparison with the dual product before its
Fourier action can be used in a geometric pairing.

The chain construction is evidence that rejecting a
naive tensor map need not reject an entire cohomological
route. It also gives a concrete next test: identify a
complex, duality and primitive projection for which the
corrected chain correspondence has the required pairing
defect. The local coefficient discrepancy specifies an
equation that such a construction must account for.

\section{Reformulated comparison}
\label{sec:reformulation}

An arithmetic transition can be tested without discarding
its correction terms. For a finite place set $S$ and support
radius $a$, suppose a specified complex has a primitive
sector $P_{S,a}$, a Hermitian pairing $I_{S,a}$, and a
complex-linear correspondence
$\Gamma_{S,a}:\T_a\to P_{S,a}$.
These are proposed geometric objects, not defined by
pulling back the Weil form. Let a transition to $(T,b)$
include an explicit map $\iota$ and a comparison of test
images. For a common test $f$, write its image difference
as $\eta f=\Gamma_{T,b}f-\iota\Gamma_{S,a}f$.

The exact pairing difference is then
\begin{align}
 &I_{T,b}(\Gamma_{T,b}f,\Gamma_{T,b}g)
       -I_{S,a}(\Gamma_{S,a}f,\Gamma_{S,a}g)\notag\\
 &\quad=E(f,g)
   +I_{T,b}(\iota\Gamma_{S,a}f,\eta g)
   +I_{T,b}(\eta f,\iota\Gamma_{S,a}g)
   +I_{T,b}(\eta f,\eta g),
 \label{eq:full-expansion}
\end{align}
where $E$ is the old-image defect defined by
\eqref{eq:old-defect} with support indices included.
An arithmetic theorem must identify the right side with
the relevant prime sum and any boundary or degree term.
Deleting the mixed terms is justified only by a proved
orthogonality identity.

For a one-prime intermediate transition, an admissible
comparison target has the form
\begin{equation}
 \Delta I(f,g)=W_p(f*g^*)+C_{S,p,a}(f,g),
 \label{eq:corrected-comparison}
\end{equation}
where $C_{S,p,a}$ is an explicitly constructed correction.
The equation is useful only if $C_{S,p,a}$ is determined
by the geometry and its behavior under further transitions
is proved. Defining $C$ as the unexplained difference of
the two sides would merely restate the problem.

At a support-complete stage the desired sign is
$I_{S,a}(x,x)\leq0$ on the appropriate primitive sector.
The full comparison must include the archimedean
distribution and the pole contribution, or reconstruct
them from an explicit degree complement. If it gives
\begin{equation}
 B_W(f,g)=-I(\Gamma f,\Gamma g)\qquad(f,g\in\T)
 \label{eq:global-target}
\end{equation}
with nonpositive $I$, then $B_W(f,f)\geq0$ on every
complex compact smooth test. The classical Weil criterion
then implies RH. Equation~\eqref{eq:global-target} and its
geometric sign are the required outcome of the program.
Neither is a consequence of equivalent Sonin norms.

The obstructions proved here impose concrete restrictions
on attempts at this outcome. Positive pointwise polar
normalization cannot preserve the entire carrier. Intrinsic
normalization must be analyzed through its compressed
operator and its test-map equations. An orthogonal
nonpositive local enlargement cannot realize the exact
intermediate $+W_p$ increment on the full pole-moment
kernel. Support-complete sign theorems and correspondences
with computed mixed terms remain compatible with these
results.

\section{Formal arithmetic statements}
\label{sec:formal}

The accompanying Lean development includes the arithmetic
coefficient calculations and the squared entire identity
obstruction. The companion sources are at
\url{https://github.com/YonedaAI/adelic-polarization-conjecture}.
The formalization uses Lean
\texttt{v4.30.0-rc2} and the pinned mathlib revision
\path{0f9072dd907c6e2e4264ab241a049cab50137f7c}.
Its theorems rely only on the usual mathlib logical
foundations, with no project arithmetic axioms or
admitted proofs. The analytic applications below remain
separate from the finite algebra they use.

The module \path{lean/APC/LocalSigns.lean} defines the
finite sum in \eqref{eq:toeplitz} for real coefficients
and the coefficient polynomial evaluation as an actual
finite sum. The theorems \path{APC.minus_exact} and
\path{APC.plus_exact} establish the two values at $11$.
Theorems with suffixes \path{minus_negative} and
\path{plus_positive} establish their strict signs.
The four radius and inverse-radius moment theorems prove
the zero polynomial evaluations for the two vectors.
The moment identities are conclusions of these proofs,
not assumptions in a sign certificate.

The module \path{lean/APC/LatticeDistribution.lean}
extends the finite calculation to a translated profile
\begin{equation}
 H_c(x)=\sum_{i,j}c_ic_jH(x-(i-j)\log p).
 \label{eq:formal-profile}
\end{equation}
Under the explicit premise that $H(k\log p)$ equals $A$
at $k=0$ and zero at every other integer $k$, its
sampling theorem identifies the prime-lattice values of
$H_c$ with $A$ times the discrete coefficient correlation.
A further theorem derives this premise from $H(0)=A$
and $H(x)=0$ for $|x|\geq\log p$. The module proves
stability of the finite prime-power sum once the cutoff
exceeds the coefficient-vector length, and its theorems
\path{APC.translated_minus_exact} and
\path{APC.translated_plus_exact} give the exact local
distribution values for the two profiles at $11$.

In Proposition~\ref{prop:polarization-obstructions:bumps},
$H=\psi*\psi^*$ and $A=\|\psi\|_2^2>0$. The Lean
profile premise does not itself construct this smooth
autocorrelation or prove its norm positive. The compact
smooth bump, convolution integral and Laplace moment
identities are proved in the paper. The all-prime packet
formula and the constrained four-bump signature are also
paper proofs. Thus the formal finite values support,
but do not replace, the complete smooth sign theorem.

The module \path{lean/APC/EntireObstruction.lean}
defines the actual entire function $\Delta_p$ and proves
its analyticity, its zero at zero and the derivative
$\log p(1-p^{-1})$ for every natural number $p>1$.
The theorem \path{APC.eulerDenominator_order_one}
deduces its exact analytic order. The final theorem
\path{APC.eulerDenominator_no_entire_square}
rules out $\Delta_pG^2=F^2$ for nonzero functions analytic
at every complex point. Finite order is deduced from
the analytic identity theorem and is not retained as
an additional premise of the entire theorem.

The passage from almost-everywhere boundary equality
to the entire squared identity is supplied by
Lemma~\ref{lem:boundary-identity}. The formal module
does not define the weighted Hilbert carrier, its
orthogonal projection or its polar decomposition.
Those operator constructions and the geometric
extension obstruction are established by the proofs
in this paper. No formal theorem asserts an arithmetic
polarization or the sign of the full Weil form.

\section{Conclusion}

The specified ambient normalization fails by a local
holomorphic parity argument, and the specified orthogonal
extension fails on exact arithmetic smooth witnesses.
The compressed Sonin normalization and the support-indexed
geometric program survive these tests with additional
comparison obligations. A successful construction must
derive its pairing, retain mixed and degree terms, and
prove the full arithmetic identity together with its sign.

\begin{thebibliography}{9}
\small
\bibitem{B04}
Jean-Fran\c{c}ois Burnol.
\emph{Two complete and minimal systems associated with the zeros of the Riemann zeta function}.
arXiv:math/0203120v7, 2004.
\url{https://arxiv.org/abs/math/0203120v7}.

\bibitem{C98}
Alain Connes.
\emph{Trace formula in noncommutative geometry and the zeros of the Riemann zeta function}.
arXiv:math/9811068v1, 1998.
\url{https://arxiv.org/abs/math/9811068v1}.

\bibitem{C26}
Alain Connes.
\emph{The Riemann Hypothesis: Past, Present and a Letter Through Time}.
arXiv:2602.04022v1, 2026.
\url{https://arxiv.org/abs/2602.04022v1}.

\bibitem{CC20}
Alain Connes and Caterina Consani.
\emph{Weil positivity and Trace formula, the archimedean place}.
arXiv:2006.13771v1, 2020.
\url{https://arxiv.org/abs/2006.13771v1}.

\bibitem{CC26}
Alain Connes and Caterina Consani.
\emph{On the Jacobian of $\overline{\operatorname{Spec}\Z}$}.
arXiv:2602.15941v1, 2026.
\url{https://arxiv.org/abs/2602.15941v1}.

\bibitem{CCM07}
Alain Connes, Caterina Consani and Matilde Marcolli.
\emph{The Weil proof and the geometry of the adeles class space}.
arXiv:math/0703392v1, 2007.
\url{https://arxiv.org/abs/math/0703392v1}.

\bibitem{CCM24}
Alain Connes, Caterina Consani and Henri Moscovici.
\emph{Zeta zeros and prolate wave operators: Semilocal adelic operators}.
arXiv:2310.18423v2, 2024.
\url{https://arxiv.org/abs/2310.18423v2}.

\bibitem{D05}
Christopher Deninger.
\emph{Arithmetic Geometry and Analysis on Foliated Spaces}.
arXiv:math/0505354v1, 2005.
\url{https://arxiv.org/abs/math/0505354v1}.

\bibitem{PartI}
Matthew Long, The YonedaAI Collaboration.
\emph{Prime-Adjoining Correspondences}.
Part I of \emph{Adelic Polarization Conjecture}, 2026.
Companion manuscript, \texttt{local-correspondence.tex}.
\end{thebibliography}
\end{document}
